\section{定位与导航}

\subsection{定位}
空中机器人在全场飞行作业，其定位需要在全局坐标系（\lstinline{map}）下连续、可靠地维护机体位姿，同时保证局部控制的实时性。我们将定位系统拆分为“全局定位 + 局部高频位姿”两级，并通过健康监测保证位姿的可用性与安全性。

\subsubsection{全局定位（激光惯性里程计）}
全局定位由 FAST-LIO 激光惯性里程计承担，融合机载 3D 激光雷达（Livox MID-360，点云输出约 200{,}000 点/秒、帧率约$\qty{10}{\Hz}$）与 IMU 数据，输出机体相对起始位姿的全局位姿。其核心优势在于：
\begin{itemize}
    \item \textbf{紧耦合高精度}：激光与惯性数据紧耦合，兼具激光在纹理环境中的高精度与 IMU 在快速运动下的高频跟踪能力；
    \item \textbf{无外置基准}：无需场地标签、磁条或 UWB 等外部基准，开机即可独立初始化，适应场地空域大、无固定信标的比赛环境；
    \item \textbf{退化鲁棒}：对弱纹理、快速旋转等退化场景具备较好的鲁棒性，配合健康监测可及时识别并规避错误定位。
\end{itemize}
里程计同时以位标志形式发布定位状态（\lstinline{/localization_status}），其中位姿可用位（bit0）与导航可用位（bit1）分别表征定位输出与建图导航的可用性，供控制与监控节点使用。

\subsubsection{局部高频位姿（IMU 插值）}
由于雷达里程计帧率有限，直接以之作为控制反馈会引入约一帧的延迟与位姿阶跃。我们利用 IMU 预积分在两帧雷达位姿之间进行插值，输出高频平滑位姿（\lstinline{/Odometry/imu_interpolation}），将有效位姿频率提升到$\qty{100}{\Hz}$量级。该高频位姿一方面作为控制回路的反馈量，另一方面经坐标变换后转发至飞控，保证飞控视觉（光流）估计器获得连续平滑的外部位姿观测。

\subsubsection{定位健康监测与门控}
在位姿转发链路中内置四级健康检查（数值有限性、坐标越界、帧间跳变、帧超时）。任一检查不通过即判定里程计不可信并锁存该状态，停止向飞控转发视觉位姿，使飞控因位姿缺失自动退出对视觉的依赖。同时将健康状态发布至 \lstinline{/odom_trustworthy} 供上层监控，并以 \lstinline{MAV_STATE_FLIGHT_TERMINATION} 状态字向飞控声明视觉定位进程异常。该机制是从源头杜绝“以错误位姿飞行”的关键安全设计。

\subsection{坐标变换与飞控数据转发}
视觉里程计与飞控使用不同的坐标系约定，需完成两次变换：
\begin{enumerate}
    \item \textbf{机体/世界约定}：MID-360 里程计采用机体 FLU、世界 ENU 约定；
    \item \textbf{飞控约定}：MAVROS 接收世界坐标系为 ENU（其内部转换至 PX4 的 NED）。
\end{enumerate}
位置按 $\boldsymbol{p}^M=(-y_L,\,x_L,\,z_L)^T$ 映射，姿态四元数按对应 $\ang{90}$ 偏航旋转叠加，速度、角速度与协方差矩阵同步旋转映射，最终发布到 \lstinline{/mavros/odometry/out}。同时转发链路支持遥控器急停开关（RC killswitch），操作手可手动切断位姿转发，实现飞行中的最后一层人工干预。

\subsection{导航}
导航层根据指令源（协议话题 / 遥控器拨杆）选择飞行模式并生成控制指令，支持航点巡航、方向机动、定点悬停与反反无人机机动。

\subsubsection{指令源与模式管理}
\begin{itemize}
    \item \textbf{话题源（Topic）}：接收电控协议下发的 \lstinline{std_msgs/UInt8MultiArray}（\lstinline{[mode, direction_flags, setpoint_index]}），按 \lstinline{setpoint_points} 航点表查表得到目标航点；
    \item \textbf{遥控器源（RC）}：通过两个三挡拨杆组合选择 6 个预设航点（主拨杆中位=悬停），实现不依赖上位机的手动航点切换；
    \item \textbf{动态来源切换}：通过来源选择拨杆在话题源与遥控器源之间运行时切换，并带超时强制回退，避免权限残留。
\end{itemize}
控制权限以 \lstinline{out_control/control_source_status}（0=不可用、1=RC、2=Topic）上报电控，供其掌握当前控制来源。

\subsubsection{航点巡航}
航点巡航采用“远距速度逼近 + 近距位置锁定”两段式策略：距离大于阈值时按比例导航以设定速度逼近目标，并在近场制动距离内线性减速；位置与航向均进入阈值后切换为位置悬停并锁定航向。整套逻辑确保飞机既能以较高速度巡航，又能在抵达航点时平稳收敛，无速度阶跃。

\subsubsection{方向机动}
方向机动按位标志输出恒定方向速度，并叠加直线保持、保高与保向修正，保证飞机沿设定方向直线飞行、高度与航向稳定。该模式常用于需要朝某一方向持续位移的场景（如初始进近、规避机动）。

\subsubsection{定点悬停}
当指令速度近似为零（小于悬停判据）或已到达目标航点时，系统切换至位置悬停模式，以当前（或目标）位置为悬停点，按位置设定点发布位置控制指令（\lstinline{type_mask} 忽略速度/加速度），并保持航向。在非 OFFBOARD 的 ALTCTL/POSCTL 模式下，悬停点以固定频率跟随当前位姿刷新，避免飞控模式切换导致的位姿跳变。

\subsubsection{反反无人机机动}
反反无人机模式以当前位置为锚点，在上下两个高度目标之间以设定频率周期切换，形成连续的垂直规避机动，用于规避对空打击。该模式的每个高度目标仍遵循航点导航逻辑与统一安全管线，且目标高度始终被工作空间高度边界钳位，保证规避机动不会越界。

\subsubsection{地理围栏}
将工作空间与最低安全高度固化为软件约束：机体靠近边界并具有向外运动趋势时，相应轴速度强制置零；高度低于最低安全高度且速度向下时，垂直速度强制置零。该约束与飞控侧的硬件级保护协同，构成双重防越界保障。

\subsection{安全保护}
导航系统从“定位、指令、权限”三个层面构建多级安全保护：
\begin{enumerate}
    \item \textbf{定位层面}：健康监测与锁存机制确保不可信位姿不会进入控制闭环与飞控；
    \item \textbf{指令层面}：统一安全管线（有限性校验、限幅、加速度约束、地理围栏）保证下发指令始终在机体能力与场地边界之内；
    \item \textbf{权限层面}：
        \begin{itemize}
            \item \textbf{输入门控}：RC 信号、里程计、定位状态、飞控状态任一超时或不可用，即停止输出所有控制指令，保持当前指令收敛；
            \item \textbf{OFFBOARD 门控}：仅在飞控处于 OFFBOARD 模式时输出运动指令，非 OFFBOARD 时发布预热悬停指令，保证模式切换安全；
            \item \textbf{一键夺权}：\lstinline{rc_stick_monitor} 在 OFFBOARD 模式下检测到操作手拨动摇杆时，主动请求飞控切换至 ALTCTL 模式，将控制权交还操作手，作为自动飞行的最终人工兜底。
        \end{itemize}
\end{enumerate}
飞行状态监测节点（\lstinline{flight_status}）将上述各环节汇总为四级飞行状态（\lstinline{/drone/flight_status}：0=与飞控断开、1=位姿不可用、2=导航不可用、3=一切正常）并上报电控，使操作手与决策层实时掌握系统健康度，实现“安全飞行，永不坠机”的总体目标。
